% !TEX root = 00_instructivo.tex
\chapter{Transductores mecánicos: acelerómetro MPU6050 y medición de altitud}

\section*{Objetivos}
\begin{itemize}
    \item Implementar un contador de pasos usando el acelerómetro MPU6050.
    \item Estimar cambios de altitud mediante doble integración de la aceleración vertical.
    \item Evaluar la precisión del sistema en caminatas reales.
\end{itemize}

\section*{Materiales y equipo}
\begin{itemize}
    \item Módulo MPU6050 (acelerómetro + giroscopio 6 DOF).
    \item Arduino Uno (o Nano) con batería portátil.
    \item Mini protoboard de 9 canales.
    \item Correa para fijar el módulo al tobillo.
    \item Cinta métrica.
\end{itemize}

\section*{Instrucciones}
\begin{enumerate}
    \item Conecte MPU6050: VCC a 5V, GND a GND, SDA a A4, SCL a A5.
    \item Calibre el acelerómetro leyendo los valores en reposo (promedio de 100 muestras).
    \item Sujete el módulo al tobillo. Programe la detección de picos en el eje Z (vertical) con un umbral dinámico.
    \item Realice 5 pruebas de 20 pasos cada una, comparando conteo manual vs conteo del sistema.
    \item Para altitud: integre dos veces la aceleración vertical (después de restar la gravedad). Compare con la altura real medida en escaleras.
    \item Grafique la aceleración en función del tiempo e identifique los picos de paso.
\end{enumerate}

\section*{Esquemático}
\begin{figure}[htbp]
\centering
\begin{circuitikz}[scale=0.8]
\draw (0,0) node[above]{Arduino Uno};
\draw (1,0.5) node[anchor=east]{5V} -- (3,0.5) node[anchor=west]{MPU6050 VCC};
\draw (1,1) node[anchor=east]{GND} -- (3,1) node[anchor=west]{GND};
\draw (1,1.5) node[anchor=east]{A4 (SDA)} -- (3,1.5);
\draw (1,2) node[anchor=east]{A5 (SCL)} -- (3,2);
\draw (4,1.75) node {I2C bus};
\end{circuitikz}
\caption{El MPU6050 se conecta directamente por I2C. Use un capacitor de 0.1µF entre 5V y GND cerca del módulo.}
\end{figure}

\section*{Preguntas de análisis}
\begin{enumerate}
    \item ¿Por qué la doble integración acumula error rápidamente? ¿Cómo se puede mitigar?
    \item ¿Qué papel juega el giroscopio en la corrección de la orientación del acelerómetro?
    \item ¿Cuál es la frecuencia de muestreo mínima recomendada para contar pasos humanos? Justifique.
    \item Proponga un filtro digital (ej. media móvil) para eliminar ruido de alta frecuencia en la señal de aceleración.
\end{enumerate}

\section*{Bibliografía}
\begin{itemize}
    \item MPU6050 Datasheet. InvenSense.
    \item “Step counting using accelerometers” – Application Note, Analog Devices.
    \item Filtro complementario para IMU – Tutorial de Arduino.
\end{itemize}
